\section{Equity}\label{sec:equity}

% ─────────────────────────────────────────────────────────────────────────────
\subsection{Forwards and dividends}

\begin{question}{\Theory \diff{1}}{Put--call parity with dividends}
Write put--call parity for an underlying paying discrete dividends.
\end{question}
\begin{answer}
$C - P = S_0 - \mathrm{PV}(\text{div}) - K\,P(0,T)$.
\end{answer}

\begin{question}{\Practice \diff{1}}{Sensitivities of the forward}
Compute by AAD the sensitivities of the 1Y forward of TICKER1 to spot, repo, dividend amounts and curve pillars, and
check them in closed form. Why are only a few pillars non-zero? \nb{equity\_underlying.ipynb}
\end{question}
\begin{answer}
Compute \texttt{F = eq.forward(T)} under a \texttt{tf.GradientTape} and take its gradient with respect to
\texttt{eq.risk\_factors()}. In closed form: $\partial F/\partial S = e^{-q_r T}/P(0,T)$, $\partial F/\partial q_r = -T F$,
$\partial F/\partial D_i = -P(0,p_i)\,e^{-q_r T}/P(0,T)$ ($p_i$ payment date) for ex-dates $t_i \le T$, 0 after. Only
the pillars around $T$ (through
$P(0,T)$) and around the payment dates of the dividends in the horizon (through their PV) move the forward: with linear
interpolation a rate depends on its two neighbouring pillars only.
\end{answer}

\begin{question}{\Theory \diff{1}}{Where is at-the-money?}
Why is the at-the-money vol of a maturity read at $K = F(T)$ rather than at $K = S_0$?
\nb{equity\_underlying.ipynb}
\end{question}
\begin{answer}
The option that is at the money for the risk-neutral distribution is the one struck at the forward (the mean of
$S_T$): its call and put have the same value. With dividends and rates $F(T)$ can be far from $S_0$ (TICKER1: 14\%
below at 5Y), so on a skewed smile $\sigma(F)$ and $\sigma(S_0)$ differ (here by up to 0.9 vol points).
\end{answer}

\begin{question}{\Theory \diff{2}}{Discrete vs continuous dividends}
How do you model discrete cash dividends and a continuous dividend yield in
Black--Scholes, and what actually changes between the two?
\nb{dividends.ipynb}
\end{question}
\begin{answer}
\textbf{Models.}
\begin{itemize}[nosep]
  \item \emph{Discrete (escrowed)}: lognormal dynamics on
        $S^* = S_0 - \mathrm{PV}(D)$, with
        $\mathrm{PV}(D) = \sum_{t_i \le T} D_i\,P(0,p_i)$ (ex-dates $t_i$, payment dates $p_i$).
  \item \emph{Continuous}: $S$ pays a yield $q$, so the stock enters as
        $S_0 e^{-qT}$. The yield matching the same forward is
        $q_{\text{eff}} = -\frac{1}{T}\ln\!\big(1 - \mathrm{PV}(D)/S_0\big)$.
\end{itemize}

\textbf{What does not change.} With $q_{\text{eff}}$ both give the same
forward $F$, and a European price depends only on $(F,\sigma,T,P(0,T))$:
\emph{same price}.

\textbf{What changes.}
\begin{itemize}[nosep]
  \item \emph{Delta}: cash dividends do not move with $S$, a yield does.
        $\partial F/\partial S$ is $e^{rT}$ (discrete) vs $F/S$ (continuous),
        so $\Delta_{\text{cont}} \approx \Delta_{\text{disc}}\,(S_0-\mathrm{PV})/S_0$.
  \item \emph{Volatility}: in the escrowed model $\sigma$ acts only on $S^*$.
        If $\sigma$ is the vol of the whole $S$ (stock drops by $D_i$ at each
        ex-date), the price is higher: until each ex-date the vol of $S^*$ is
        $\sigma\,S/(S-\mathrm{PV}_t)$, so
        $\sigma^{*2} T \approx \sigma^2 \sum_k \big(S_0/(S_0-\mathrm{PV}_k)\big)^2 \Delta t_k$
        (the flat $\sigma\,S_0/(S_0-\mathrm{PV})$ overshoots when the dividends are late).
  \item \emph{American options}: early exercise of a call is optimal just before
        an ex-date; a yield smears the dividends over time and misses
        this timing.
  \item \emph{Long maturities}: yields are proportional and never push $S$
        below zero; cash dividends can be large relative to a low $S$.
\end{itemize}
\end{answer}

\begin{question}{\Practice \diff{2}}{Sticky yield: AAD delta vs bump}
Price a 1Y call with \texttt{BlackScholesPricer(dividend\_model="continuous")} on an underlying with discrete
dividends. Compare the AAD delta with a bump-and-reprice of the spot. Why do they differ?
\nb{dividends.ipynb}
\end{question}
\begin{answer}
The pricer computes $q_{\text{eff}} = -\ln(1 - \mathrm{PV}/S)/T$ with $S$ under \texttt{stop\_gradient}: the AAD delta
holds the yield (sticky yield, $\Delta_{\text{disc}}(S_0-\mathrm{PV})/S_0$). A bump moves $S$ in $q_{\text{eff}}$ too, the
cash amounts stay fixed: it gives the discrete (sticky cash) delta. Same price, two conventions.
\end{answer}

\begin{question}{\Theory \diff{2}}{Early exercise and the dividend model}
In the escrowed model the American call is worth the European one; with the equivalent yield it has a premium. Which
is right for a stock paying a cash dividend?
\nb{dividends.ipynb}
\end{question}
\begin{answer}
Neither exactly. A call is exercised early only just before an ex-date, to capture the dividend. The escrowed model
has carry $r - q_{\text{repo}} \ge 0$ and no drop at all, so it never exercises; the yield spreads the dividend over
time and exercises at the wrong moments. The real premium needs the cash drop in the dynamics (tree or PDE with jumps
of the stock at the ex-dates). For puts it is the opposite: the escrowed model has the premium (carry $\approx r$), the
yield model almost none.
\end{answer}

\begin{question}{\Practice \diff{2}}{Checking an American approximation}
How do you check the Bjerksund--Stensland price of \texttt{BlackScholesPricer} on American options?
\end{question}
\begin{answer}
Against a binomial tree (CRR, a few thousand steps) on the same inputs read from the priced product:
\texttt{spot\_net}, \texttt{risk\_free\_rate}, \texttt{volatility}, \texttt{time\_to\_maturity} and the carry
$q = r - \ln(F/S_{\text{net}})/T$ from \texttt{forward}. The approximation is a lower bound within about 1\%; the
check found a wrong trigger price in the implementation (early-exercise premium 3--4 times too small).
\end{answer}

% ─────────────────────────────────────────────────────────────────────────────
\subsection{Options and Black--Scholes}

\begin{question}{\Theory \diff{1}}{The Black--Scholes model}
Define the Black--Scholes model for an equity, state its main assumptions and
outline how the European call pricing formula is obtained.
\end{question}
\begin{answer}
\textbf{Model.} The stock follows a geometric Brownian motion; under the
risk-neutral measure $\Q$, with continuous dividend yield $q$:
\[
  \dd S_t = (r - q)\,S_t \dd t + \sigma S_t \dd W_t^{\Q}
  \quad\Rightarrow\quad
  S_T = S_0\, e^{(r - q - \frac{1}{2}\sigma^2)T + \sigma\sqrt{T}\,Z},\; Z\sim\Nc(0,1).
\]

\textbf{Assumptions.}
\begin{itemize}[nosep]
  \item constant volatility $\sigma$, rate $r$ and dividend yield $q$;
  \item lognormal prices, continuous paths (no jumps);
  \item frictionless market: continuous trading, no transaction costs or
        taxes, short selling allowed, divisible assets;
  \item no arbitrage.
\end{itemize}

\textbf{Derivation.} Two equivalent routes.
\begin{enumerate}[nosep]
  \item \emph{Hedging (PDE).} Hold the option and short $\Delta$ shares:
        $\Pi = V - \Delta S$. By It\^o, choosing $\Delta = \partial_S V$
        cancels the $\dd W$ term, so $\Pi$ is riskless and must earn $r$:
        \[
          \partial_t V + (r - q) S\,\partial_S V
          + \tfrac{1}{2}\sigma^2 S^2\,\partial_{SS} V - rV = 0,
          \qquad V(T,S) = (S - K)^+ .
        \]
  \item \emph{Risk-neutral pricing.} Solving the PDE (Feynman--Kac) is
        equivalent to
        $C = e^{-rT}\,\E^{\Q}\big[(S_T - K)^+\big]$, which is computed
        in closed form using the lognormal $S_T$.
\end{enumerate}

\textbf{Result.}
\[
  C = S_0 e^{-qT} \Nc(d_1) - K e^{-rT} \Nc(d_2),
  \qquad
  d_{1,2} = \frac{\ln(S_0/K) + (r - q \pm \frac{1}{2}\sigma^2)T}{\sigma\sqrt{T}} .
\]

\textbf{Intuition.} The drift $\mu$ of the real-world dynamics does not appear:
the risk is hedged away, so only $\sigma$ matters.
$\Nc(d_2) = \Q(S_T > K)$ is the risk-neutral exercise probability;
$S_0 e^{-qT}\Nc(d_1)$ is the value of receiving the stock if exercised, and
$\Nc(d_1)$ (times $e^{-qT}$) is the delta.
\end{answer}

\begin{question}{\Theory \diff{2}}{Gamma close to expiry}
What happens to gamma and theta of an at-the-money option as $T \to 0$, and of an out-of-the-money one?
\nb{black\_scholes.ipynb}
\end{question}
\begin{answer}
ATM, $\Gamma \approx \frac{n(d_1)}{S\sigma\sqrt T} \to \infty$ as $1/\sqrt T$ (the payoff kink at $K$), and theta, its mirror in the
PDE ($\theta \approx -\tfrac12 \sigma^2 S^2 \Gamma$), explodes too. OTM (or ITM) both go to zero: $d_1 \to \pm\infty$ and
$n(d_1)$ vanishes faster than $1/\sqrt T$ grows. Near expiry the risk concentrates on the strike.
\end{answer}

\begin{question}{\Theory \diff{2}}{Gamma trading}
You buy a 1Y ATM call at 30\% implied vol and delta-hedge it daily. What is your P\&L if the stock realises 25\%,
30\%, 35\%? What are you betting on? \nb{vanilla\_options.ipynb}
\end{question}
\begin{answer}
$\text{P\&L} \approx \tfrac12\int \Gamma S^2(\sigma_r^2 - \sigma_i^2)\,\dd t \approx \mathcal{V}(\sigma_r - \sigma_i)$: about $-4.5$, $0$, $+4.5$
on a 27 premium (vega 0.90 per vol point). Long gamma is long realised vol against the implied: you pay theta at
$\sigma_i$ and earn gamma on the realised moves. The P\&L is path dependent (vol near the strike counts most).
\end{answer}

\begin{question}{\Practice \diff{2}}{Hedging error and rebalancing}
With realised = implied vol, how does the dispersion of the delta-hedged P\&L depend on the number of rebalances?
\nb{vanilla\_options.ipynb}
\end{question}
\begin{answer}
The mean stays at zero (the hedge replicates the option); the standard deviation falls as $1/\sqrt n$ (TICKER1:
6.5, 3.2, 1.5, 0.73 for 12, 52, 252, 1000 rebalances): halving the error costs four times the trades, which in practice
sets the hedging frequency against transaction costs.
\end{answer}

\begin{question}{\Theory \diff{2}}{Digital options and the skew}
Why is a digital call worth more on a market with negative skew than its Black--Scholes price at the vol of the
strike? How do you price it consistently with the smile? \nb{digital\_options.ipynb}
\end{question}
\begin{answer}
A digital is minus the strike derivative of the call: $D = -\partial C/\partial K = P(0,T)N(d_2) - \mathcal{V}\,\partial\sigma/\partial K$.
With $\partial\sigma/\partial K < 0$ the second term adds value (TICKER1 1Y ATM: $+0.07$ on 0.43). Consistent price: a tight
call spread on the market vols, or a model that reprices the smile (local vol); it depends on how the smile is
interpolated between the quoted strikes, because it is its slope that matters.
\end{answer}

\begin{question}{\Theory \diff{2}}{The vega flip of a digital}
Why does the vega of a long digital call change sign, and how do you hedge a short digital near expiry?
\nb{digital\_options.ipynb}
\end{question}
\begin{answer}
$\mathcal{V} = -P(0,T)\,n(d_2)\,d_1/\sigma$: out of the money ($d_1 < 0$) more vol raises the chance of finishing above $K$,
in the money it raises the chance of falling below; the sign flips at $d_1 = 0$, near the strike, like the gamma.
Near expiry the delta is a spike of height $\sim 1/(S\sigma\sqrt T)$ at the strike: delta-hedging the digital alone
means huge trades on small moves (pin risk). Hedge statically with the over-replicating call spread
(long $1/h$ calls at $K-h$, short $1/h$ at $K$): it matches delta, gamma and vega with bounded greeks and pays at least
the digital.
\end{answer}

\begin{question}{\Practice \diff{2}}{Monte Carlo greeks of a digital}
Price a digital call by Monte Carlo and take its AAD delta, first with the hard indicator
(\texttt{digital\_smoothing} $= 0$), then with a smoothing of $0.01$. Why does the first give nothing? \nb{digital\_options.ipynb}
\end{question}
\begin{answer}
The hard payoff is an indicator: on each path its derivative w.r.t.\ the spot is zero (and TF has no gradient
through the boolean), while the price moves by the probability mass crossing $K$. A sigmoid of width $\varepsilon K$
makes the payoff smooth: the pathwise delta becomes the Black--Scholes one up to noise, at the cost of a small
price bias (second order in $\varepsilon$). Alternatives: likelihood-ratio weights, or a call-spread payoff.
\end{answer}

\begin{question}{\Practice \diff{2}}{European vs American}
For TICKER1 at 1Y, which options have an early-exercise premium in the escrowed and in the equivalent-yield model,
and how does it move with the maturity? \nb{black\_scholes.ipynb}
\end{question}
\begin{answer}
Escrowed (carry $r - q_r > 0$): calls none, puts yes (1.4\% of the ATM put, more for ITM strikes). Equivalent yield
($q_{\text{eff}} > r$): calls yes (2.6\% ATM, more ITM), puts almost none. The premium grows with $T$ and jumps at each
ex-date that enters the option life (lower escrowed spot, higher equivalent yield). Below the exercise boundary the
American put is worth $K - S$: delta $-1$.
\end{answer}

% ─────────────────────────────────────────────────────────────────────────────
\subsection{Volatility: smile, local vol, displaced diffusion}

\begin{question}{\Practice \diff{1}}{Which nodes carry the vega?}
A price uses the vol surface at one $(K, T)$ between the nodes. Which nodes of \texttt{volatility\_matrix} does it
depend on? Check it by bumping one node with \texttt{assign}. \nb{market\_handles.ipynb}
\end{question}
\begin{answer}
\texttt{volatility(K, T)} is bilinear and flat outside: only the 4 nodes around $(K, T)$ (2 on an edge, 1 at a
corner), with the bilinear weights. Hence the bucketed vega (gradient w.r.t.\ the matrix) of a Black--Scholes price
sits on those nodes; a model that builds its own surface (local vol) spreads it further.
\end{answer}

\begin{question}{\Theory \diff{2}}{Dupire formula}
% TODO
\end{question}
\begin{answer}
% TODO
\end{answer}

% ─────────────────────────────────────────────────────────────────────────────
\subsection{Monte Carlo and AAD}

\begin{question}{\Theory \diff{1}}{Monte Carlo standard error}
How does the error of a Monte Carlo estimator scale with the number of paths $N$?
\end{question}
\begin{answer}
$\mathrm{SE} = \hat\sigma/\sqrt{N}$: halving the error requires $4\times$ the paths.
\end{answer}

\begin{question}{\Practice \diff{2}}{Greeks via automatic differentiation}
Price a 1Y call with \texttt{VanillaMCPricer} and get delta, vega per node, rho per pillar and the dividend
sensitivities in one go. Compare with Black--Scholes. \nb{monte\_carlo\_gbm.ipynb}
\end{question}
\begin{answer}
Price with \texttt{autodiff=True}, then take \texttt{pricer.tape.gradient} of the price with respect to
\texttt{eq.risk\_factors()}: a dict keyed by market key with the gradient of every variable (spot, vol matrix, curve
pillars, dividend amounts). One backward pass costs a few prices, whatever the
number of greeks. With 200k paths the MC greeks match Black--Scholes to the MC noise (delta 0.5591 vs 0.5593).
\end{answer}

\begin{question}{\Theory \diff{2}}{Common random numbers}
Why does a bump-and-reprice delta of a Monte Carlo price need the same random numbers in the two prices?
\nb{monte\_carlo\_gbm.ipynb}
\end{question}
\begin{answer}
$(V(S+h) - V(S-h))/2h$: with independent draws each price carries noise of size s.e., so the estimate has noise
$\sim \text{s.e.}/h$, exploding as $h \to 0$ (0.1 bump: 1.41 instead of 0.56). With the same draws the noise is common and
cancels: the difference is the change of the payoff on the same paths, close to the pathwise derivative.
\end{answer}

\begin{question}{\Practice \diff{2}}{Variance reduction}
Reduce the standard error of a Monte Carlo call with antithetic variates and with a control variate. Which works
better here, and why? \nb{monte\_carlo\_gbm.ipynb}
\end{question}
\begin{answer}
Antithetic ($Z$, $-Z$) helps where the payoff is close to linear in $Z$: on an ATM call 0.151 $\to$ 0.124. The control
variate $S_T$ (mean = forward, known exactly) is highly correlated with the call payoff: 0.151 $\to$ 0.064, a factor
of about 5 in paths. The best control is the most correlated quantity with a known expectation.
\end{answer}
